%==============================================================
% DOCUMENTCLASS
%==============================================================
\documentclass[12pt,oneside]{book}

%==============================================================
% METADATOS
%==============================================================
\newcommand{\universidad}{Universidad de Buenos Aires}
\newcommand{\facultad}{Facultad de Derecho}
\newcommand{\materia}{Inglés Jurídico}
\newcommand{\titulo}{Apuntes de Inglés Jurídico}
\newcommand{\subtitulo}{Lectocomprensión y terminología jurídica angloamericana}
\newcommand{\autor}{Isaac Edgar Camacho Ocampo}
\newcommand{\anio}{2026}

%==============================================================
% PAQUETES
%==============================================================
\usepackage[utf8]{inputenc}
\usepackage[T1]{fontenc}
% Babel no está disponible en esta instalación; se omite para evitar fallos de compilación.

\usepackage[
    a4paper,
    top=2.5cm,
    bottom=2.5cm,
    left=3cm,
    right=2.5cm,
    headheight=15pt
]{geometry}

\usepackage{amsmath}
\usepackage{amssymb}
\usepackage{mathtools}

\usepackage{graphicx}
\usepackage{float}
\usepackage{caption}
\usepackage{subcaption}

\usepackage{booktabs}
\usepackage{array}
\usepackage{longtable}
\usepackage{tabularx}

\usepackage{enumitem}

\usepackage{xcolor}
\usepackage[most]{tcolorbox}

\usepackage{fancyhdr}
\usepackage{ragged2e}

\usepackage[
    colorlinks=true,
    linkcolor=blue!50!black,
    citecolor=green!40!black,
    urlcolor=blue!60!black,
    pdftitle={\titulo},
    pdfauthor={\autor},
    pdfsubject={Apuntes universitarios},
    bookmarks=true
]{hyperref}

%==============================================================
% CONFIGURACIÓN
%==============================================================
\graphicspath{
    {./img/}
    {./graficos/}
    {./figuras/}
}

\setcounter{tocdepth}{2}

\setlength{\parindent}{0pt}
\setlength{\parskip}{0.5em}

\setlist[itemize]{
    topsep=0.3em,
    itemsep=0.2em
}

\setlist[enumerate]{
    topsep=0.3em,
    itemsep=0.2em
}

\clubpenalty=10000
\widowpenalty=10000

\renewcommand{\arraystretch}{1.2}

\newcolumntype{L}{>{\raggedright\arraybackslash}X}
\newcolumntype{P}[1]{>{\raggedright\arraybackslash}p{#1}}

%==============================================================
% COMANDOS
%==============================================================
\newcommand{\abs}[1]{\left|#1\right|}
\newcommand{\norm}[1]{\left\lVert#1\right\rVert}
\newcommand{\vect}[1]{\mathbf{#1}}
\newcommand{\deriv}[2]{\frac{d#1}{d#2}}
\newcommand{\pderiv}[2]{\frac{\partial#1}{\partial#2}}

%==============================================================
% ENTORNOS / CAJAS
%==============================================================
\newtcolorbox{definition}[1]{
    enhanced,
    breakable,
    title={Definición: #1},
    fonttitle=\bfseries,
    colback=blue!3,
    colframe=blue!50!black,
    boxrule=0.8pt,
    arc=2mm
}

\newtcolorbox{important}{
    enhanced,
    breakable,
    title={Importante},
    fonttitle=\bfseries,
    colback=yellow!8,
    colframe=orange!70!black,
    boxrule=0.8pt,
    arc=2mm
}

\newtcolorbox{example}[1]{
    enhanced,
    breakable,
    title={Ejemplo: #1},
    fonttitle=\bfseries,
    colback=green!4,
    colframe=green!40!black,
    boxrule=0.8pt,
    arc=2mm
}

\newtcolorbox{formula}[1]{
    enhanced,
    breakable,
    title={Fórmula / Regla: #1},
    fonttitle=\bfseries,
    colback=purple!4,
    colframe=purple!50!black,
    boxrule=0.8pt,
    arc=2mm
}

\newtcolorbox{exercise}[1]{
    enhanced,
    breakable,
    title={Ejercicio: #1},
    fonttitle=\bfseries,
    colback=gray!5,
    colframe=gray!60!black,
    boxrule=0.8pt,
    arc=2mm
}

\newtcolorbox{note}{
    enhanced,
    breakable,
    title={Nota},
    fonttitle=\bfseries,
    colback=white,
    colframe=gray!50,
    boxrule=0.6pt,
    arc=2mm
}

%==============================================================
% CONFIGURACIÓN FINAL (ENCABEZADOS Y PIES)
%==============================================================
\renewcommand{\chaptermark}[1]{%
    \markboth{\ifnum\value{chapter}>0 \thechapter.\ \fi#1}{}%
}

\pagestyle{fancy}
\fancyhf{}

\fancyhead[L]{\nouppercase{\leftmark}}
\fancyhead[R]{\materia}

\fancyfoot[C]{\thepage}

\renewcommand{\headrulewidth}{0.4pt}
\renewcommand{\footrulewidth}{0pt}

\fancypagestyle{plain}{
    \fancyhf{}
    \fancyfoot[C]{\thepage}
    \renewcommand{\headrulewidth}{0pt}
}

%==============================================================
% DOCUMENTO
%==============================================================
\begin{document}

%--------------------------------------------------------------
% PORTADA
%--------------------------------------------------------------
\begin{titlepage}

\thispagestyle{empty}

\begin{center}

\vspace*{1cm}

{\large \textsc{\universidad}}

\vspace{0.2cm}

{\large \textsc{\facultad}}

\vspace{3cm}

{\Huge \textbf{\materia}}

\vspace{1cm}

{\LARGE \titulo}

\vspace{0.4cm}

{\large \subtitulo}

\vspace{0.8cm}

\textit{\small Comprender es el primer paso para convertir el estudio en conocimiento.}

\vspace{1.2cm}

\rule{0.70\textwidth}{0.5pt}

\vspace{0.5cm}

{\large Apuntes de estudio}

\vspace{0.5cm}

\rule{0.70\textwidth}{0.5pt}

\vfill

{\large \autor}

\vspace{0.3cm}

{\large \anio}

\end{center}

\vfill

\begin{flushleft}
\textit{Este es un modesto aporte a los estudiantes de ``Inglés Jurídico'' no pretende sustituir el material oficial de la materia.}
\end{flushleft}

\end{titlepage}

%--------------------------------------------------------------
% ÍNDICE
%--------------------------------------------------------------
\frontmatter

\tableofcontents

%--------------------------------------------------------------
% PRESENTACIÓN DE LA CLASE
%--------------------------------------------------------------
\chapter{Presentación de la clase}

\section*{Datos de la clase}

\begin{tabularx}{\textwidth}{@{}P{3.2cm}L@{}}
\toprule
\textbf{Materia} & Inglés Jurídico (\textit{Legal English}) \\
\textbf{Institución} & Facultad de Derecho, Universidad de Buenos Aires \\
\textbf{Fecha} & 26 de octubre de 2026 \\
\textbf{Clase} & Repaso previo al examen: texto sobre el servicio de jurado, procesos civil y penal angloamericanos, contratos y vocabulario técnico \\
\bottomrule
\end{tabularx}

\section*{Temas tratados}

La clase funcionó como repaso integral antes del examen. En primer lugar se aclaró la \textbf{modalidad} de la evaluación: una única pregunta de redacción, respuestas siempre en español, traducción de los términos que aparezcan en inglés, uso de diccionarios jurídico y común, y valor del \textbf{sentido común} para leer un texto.

Después se recorrieron los grandes temas del curso: el \textbf{proceso penal} (\textit{arraignment}, \textit{plea bargaining}, \textit{nolo contendere}, \textit{indictment} y estándares de prueba), la \textbf{Constitución británica y los contratos} (\textit{offer and acceptance}, \textit{unilateral contract} y \textit{bilateral contract}), el \textbf{jurado} (\textit{excuses}, \textit{challenges}, \textit{voir dire}, \textit{jury panel} y \textit{juror}), las \textbf{\textit{motions}} y los \textbf{\textit{equitable remedies}}, la etapa de \textbf{\textit{discovery}} y las \textbf{\textit{depositions}}, el \textbf{\textit{trier of fact}} y el término \textbf{\textit{duress}}.

\section*{Utilidad de los contenidos}

\subsection*{¿Para qué sirve esto para la vida profesional?}

Sirve para trabajar como abogado o abogada en un mundo donde los contratos, los litigios y la doctrina cruzan fronteras. Leer jurisprudencia o contratos en inglés exige distinguir \textbf{falsos amigos} (\textit{equitable} no es «equitativo»; \textit{motion} no es «moción») y saber \textbf{a qué institución propia equivale} cada figura (\textit{discovery} equivale aproximadamente a una etapa probatoria preliminar; \textit{information}, a la elevación a juicio). Eso permite redactar, traducir y asesorar a clientes o colegas extranjeros sin errores que cambian el sentido jurídico de un documento, y comprender qué puede pasar en un juicio con jurado de otro sistema.

\subsection*{¿Es útil para la vida en general?}

Sí. Más allá de la profesión, la clase entrena habilidades de uso permanente: \textbf{leer con sentido común} y no quedarse trabado en una palabra suelta, \textbf{interpretar el contexto} antes de buscar en el diccionario, \textbf{comparar sistemas} para entender el propio y \textbf{estudiar con criterio} (llevar lo justo, priorizar los apuntes). Además, entender cómo funciona un jurado, un acuerdo con el fiscal o un intercambio de pruebas ayuda a comprender las noticias, series y casos de otros países, y a razonar con más claridad ante cualquier decisión.

\section*{Hilo conductor: el collar de Lucía}

Una historia cotidiana atraviesa todos los bloques de la clase. Lucía pierde su collar y cuelga carteles: «recompensa para quien lo encuentre».

\begin{table}[H]
\centering
\begin{tabularx}{\textwidth}{P{3.6cm}L}
\toprule
\textbf{Tema} & \textbf{Qué pasa en la historia de Lucía} \\
\midrule
El examen &
Lucía redacta el cartel de la recompensa y quien lo lee necesita un diccionario jurídico, uno común (\textit{necklace}) y, sobre todo, sentido común para interpretarlo. \\
Proceso penal &
Si alguien se queda con el collar y es acusado, le leen los cargos (\textit{arraignment}) y le preguntan cómo se declara; el jurado necesita convicción más allá de toda duda razonable. \\
Constitución y contratos &
El cartel es una oferta a persona indeterminada: quien encuentra el collar acepta cumpliendo la condición (\textit{unilateral contract}). \\
Jurado &
Si hay juicio con jurado, algunas personas se excusan y las partes recusan con o sin causa (\textit{excuses} y \textit{challenges}). \\
\textit{Motions} y \textit{remedies} &
Las partes hacen pedidos al juez (\textit{motions}) y, si Lucía no paga, se puede pedir que cumpla lo prometido (\textit{specific performance}). \\
\textit{Indictment} y \textit{discovery} &
Antes del juicio, las partes intercambian pruebas, toman \textit{depositions} y deciden si hay acuerdo; el \textit{indictment} o la \textit{information} habilita seguir adelante. \\
\textit{Plea} y selección &
Entre negociaciones con el fiscal (\textit{plea bargaining}) y la selección del jurado (\textit{voir dire}), se forma el \textit{jury panel} y luego el \textit{jury}. \\
\textit{Trier of fact} y cierre &
El jurado resuelve los hechos y el juez el derecho (\textit{trier of fact}); si hubo presión indebida, se habla de \textit{duress}. \\
\bottomrule
\end{tabularx}
\end{table}

\mainmatter

%==============================================================
% CAPÍTULO 1
%==============================================================
\chapter[El examen]{El examen: modalidad, consignas y materiales}

\section{Modalidad y duración}

El examen está previsto para resolverse en un máximo de \textbf{dos horas}. Contiene una \textbf{única pregunta de redacción}; el resto de las consignas se completa en el propio examen: hay preguntas de verdadero o falso, de \textit{multiple choice} y de otras modalidades de completado.

El examen fue armado de modo que no haya que escribir mucho. La razón expresada es que, al corregir respuestas redactadas, resulta difícil distinguir si el estudiante no comprendió el tema o si lo comprendió pero no supo expresarlo, de modo que la corrección termina siendo una conjetura.

\section{Idioma de respuesta y términos de la consigna}

\textbf{Todas las respuestas se redactan en español}; no hay que escribir nada en inglés.

Si una consigna incluye un término en inglés, ello es \textbf{deliberado}: se espera que el estudiante lo traduzca o dé su equivalente en español, en lugar de repetir el término extranjero.

\begin{example}{La etapa de \textit{discovery} en la consigna}
Si la pregunta menciona «\textit{discovery}» (por ejemplo, «¿qué es la etapa de \textit{discovery}?»), el estudiante debe responder con la traducción o el equivalente en español (por ejemplo, la etapa de intercambio de información probatoria), sin repetir el término en inglés.
\end{example}

\section{Materiales permitidos y sentido común}

\subsection{Qué materiales llevar}

Se puede llevar el material que cada uno desee (papel y material impreso), pero \textbf{no se permiten dispositivos}. Conviene llevar hojas propias, porque quien tiene letra grande puede no alcanzar con la hoja impresa de la consigna.

La recomendación es no llevar una cantidad excesiva de material, sobre todo para quien no tiene un buen nivel de inglés: llevar «la biblioteca» da sensación de resguardo, pero se pierde mucho tiempo buscando. Lo más útil son las \textbf{anotaciones de clase}.

\subsection{El sentido común como herramienta de lectura}

\begin{important}
Hay que leer realmente el texto y no quedarse con el dato anecdótico del diccionario o del glosario: \textbf{interpretar el contexto}, porque muchas veces la respuesta está «delante de la nariz». Es más importante interpretar que encontrar la palabra precisa.
\end{important}

Quienes cursan ya conocen mucho derecho y analizaron numerosas cuestiones del sistema extranjero; por eso, con frecuencia, se trata de interpretar más que de hallar el término exacto. Llevar material «por las dudas» puede jugar en contra.

\section{Diccionarios}

Se necesitan \textbf{dos tipos de diccionario}, que pueden repartirse entre quienes rinden: uno trae uno y otro trae el otro.

\begin{table}[H]
\centering
\begin{tabularx}{\textwidth}{P{3.8cm}LL}
\toprule
\textbf{Tipo de diccionario} & \textbf{Para qué sirve} & \textbf{Ejemplo mencionado} \\
\midrule
Diccionario jurídico & Términos técnicos del derecho & El de Cabanellas, disponible en la biblioteca (sector de referencia) \\
Diccionario bilingüe común & Vocabulario general que aparece en el texto & \textit{Necklace} (collar) \\
\bottomrule
\end{tabularx}
\end{table}

\begin{example}{El collar perdido (\textit{necklace})}
En un texto de contratos visto en clase aparecía el ejemplo de un collar perdido. Aunque \textit{necklace} no es un término jurídico, quien no lo conoce no entiende el caso: por eso hace falta, además del diccionario jurídico, un diccionario bilingüe común.
\end{example}

Los diccionarios pueden compartirse, pero conviene que circulen varios ejemplares y que no se los estén pasando de mano en mano, sobre todo porque el aula donde se rinde es grande. En la biblioteca hay varios ejemplares del Cabanellas, y recientemente se compraron más porque los estudiantes lo buscan mucho.

\begin{note}
En la biblioteca existen dos sectores. \textbf{Préstamos} (en el primer piso) es el sector donde se pueden retirar libros de cualquier tema para llevar a casa. El diccionario se encuentra en \textbf{referencia}: puede retirarse, pero debe solicitarse desde ese sector.
\end{note}

\section{La pregunta de desarrollo}

Todas las preguntas del examen se responden con el texto, salvo la \textbf{última}, que es la de desarrollo. Esta también está relacionada con el texto, pero apunta a que se conteste con \textbf{contenido teórico trabajado en clase}, por ejemplo, alguna diferencia entre negocios jurídicos o algún tema particular abordado en el curso. Es la única pregunta cuya respuesta no se desprende del texto.

\section{Qué se evalúa}

En la pregunta de desarrollo no interesa tanto la redacción como el uso de los \textbf{términos precisos} (\textit{discovery}, \textit{peremptory challenge} y otros trabajados en clase). Se trata de términos técnicos, y la terminología jurídica es parte de lo que la materia se propone abordar. Por eso no es indistinto escribir «moción» o «pedido»: hay que traducir correctamente.

% TODO: contenido incompleto o ambiguo en las notas originales (comentario de un estudiante sobre «sentence» como «sentencia» y pasaje poco claro en el audio)

\section{Cuadro Cornell: el examen}

\begin{table}[H]
\centering
\begin{tabularx}{\textwidth}{
    >{\raggedright\arraybackslash}p{0.28\textwidth}
    >{\raggedright\arraybackslash}X
}
\toprule
\textbf{Preguntas / palabras clave} &
\textbf{Notas / respuestas} \\
\midrule

\textbf{¿Cómo está armado el examen?} &
Dura como máximo dos horas. Tiene una única pregunta de redacción; el resto se completa (verdadero o falso, \textit{multiple choice}). \\
\midrule[0.3pt]

\textbf{¿En qué idioma se responde y qué se hace con los términos en inglés de la consigna?} &
Todo en español. Si la consigna trae un término en inglés (por ejemplo, \textit{discovery}), es a propósito: hay que traducirlo o dar su equivalente, no repetirlo. \\
\midrule[0.3pt]

\textbf{¿Qué materiales conviene llevar y qué papel cumple el sentido común?} &
Hojas y el material que cada uno quiera, sin dispositivos, pero sin llenarse de cosas porque se pierde tiempo buscando; lo más útil son las anotaciones de clase. Hay que leer el texto de verdad e interpretar el contexto: a veces la respuesta está «delante de la nariz». \\
\midrule[0.3pt]

\textbf{¿Qué diccionarios hacen falta?} &
Uno jurídico (por ejemplo, el de Cabanellas, en el sector de referencia de la biblioteca) y uno bilingüe común (\textit{necklace} no es un término jurídico). Pueden compartirse. \\
\midrule[0.3pt]

\textbf{¿Qué se espera de la pregunta de desarrollo?} &
Está vinculada al texto, pero se responde con contenido teórico visto en clase (por ejemplo, una diferencia entre figuras jurídicas). Interesan más los términos precisos que la redacción. \\

\midrule

\multicolumn{2}{p{\dimexpr\textwidth-2\tabcolsep\relax}}{
\textbf{Resumen:}
El examen combina una pregunta de redacción con consignas de completado, se responde siempre en español y exige traducir los términos en inglés de la consigna. Importan la interpretación del contexto y el uso preciso de la terminología; los diccionarios (jurídico y común) y las anotaciones de clase son apoyo, no sustituto del estudio.
} \\

\bottomrule
\end{tabularx}
\end{table}

%==============================================================
% CAPÍTULO 2
%==============================================================
\chapter[Proceso penal angloamericano]{Proceso penal angloamericano: arraignment, plea y estándares de prueba}

\section{Arraignment: la audiencia en la que el acusado se declara}

Antes del \textit{arraignment} el acusado pasó por la policía, la detención y el arresto. El \textit{arraignment} es la \textbf{etapa} en que se le informa formalmente de qué se lo acusa y se le pregunta cómo se declara. Es una etapa bastante preliminar, parecida a nuestra instrucción, y se ubica \textbf{antes del acuerdo con el fiscal}.

\begin{definition}{\textit{Arraignment} (audiencia de lectura de cargos)}
Audiencia inicial del proceso penal, posterior a la detención y el arresto, en la que se informa formalmente al acusado de qué se lo acusa (lectura formal de cargos) y, a continuación, en esa misma audiencia, se le pregunta cómo se declara.
\end{definition}

Si en ese momento el acusado dice «me declaro culpable», lo hace a cambio de algún beneficio que le otorgó el fiscal; de lo contrario, no habría razón para declararse culpable. Todo ocurre en esa misma audiencia.

\subsection{Modificación de la fianza}

En esa primera audiencia formal de acusación o de instrucción (la lectura de los cargos y la decisión de avanzar), el juez puede \textbf{ajustar el monto de la fianza}, en función de los indicios probatorios presentados o del comportamiento del propio imputado.

% TODO: contenido incompleto o ambiguo en las notas originales (la transcripción original decía «amianza inicial»; se interpretó como «fianza» porque el contexto habla de ajustar su monto)

\section{Plea bargaining: la negociación previa con el fiscal}

El \textit{plea bargaining} es un \textbf{proceso previo}, una negociación con el fiscal. Si al fiscal le conviene que la persona se declare culpable (por ejemplo, porque eso le permitirá obtener información sobre otros hechos o sobre otros sospechosos, en una causa grande con más partes involucradas), le propone algo: declararse culpable a cambio de un beneficio, como la reducción de la pena.

Cuando en la audiencia de \textit{arraignment} se le pregunta cómo se declara, el acusado ya acordó previamente con el fiscal. Ante el juez debe decir \textit{guilty}, \textit{not guilty} o \textit{nolo contendere}.

\begin{table}[H]
\centering
\begin{tabularx}{\textwidth}{LL}
\toprule
\textbf{\textit{Plea bargaining}} & \textbf{\textit{Arraignment}} \\
\midrule
Proceso previo: negociación con el fiscal. & Audiencia en la que se leen los cargos y se pregunta cómo se declara. \\
El fiscal puede ofrecer reducción de pena a cambio de que el acusado se declare culpable (por ejemplo, para obtener información sobre otros sospechosos). & El acusado declara ante el juez: \textit{guilty}, \textit{not guilty} o \textit{nolo contendere}. \\
\bottomrule
\end{tabularx}
\end{table}

Si lo que ofrece el fiscal no sirve o no interesa, o si la persona considera que puede demostrar su inocencia, dirá \textit{not guilty}. Si, en cambio, la situación del acusado es muy comprometida y el fiscal le ofreció un beneficio, probablemente diga \textit{guilty}: se ahorra tiempo, porque no se avanza a la etapa probatoria, y además recibe un beneficio, como la reducción de la pena.

\section{Las formas de declararse: guilty, not guilty y nolo contendere}

\begin{table}[H]
\centering
\begin{tabularx}{\textwidth}{P{3.6cm}L}
\toprule
\textbf{Forma de declararse} & \textbf{Significado} \\
\midrule
\textit{Guilty} & Culpable. \\
\textit{Not guilty} & No culpable. \\
\textit{Nolo contendere} & «No voy a contradecir, no voy a confrontar, no voy a objetar»: ni blanco ni negro. Funciona como una admisión de culpabilidad, pero inhibe la instancia civil por el mismo hecho. \\
\bottomrule
\end{tabularx}
\end{table}

Ante la objeción de un estudiante (quien tiene todas las de ganar diciendo \textit{not guilty}, ¿por qué se autoincriminaría?), se explicó que el \textit{nolo contendere} funciona en definitiva como una admisión de culpabilidad, pero su principal efecto es que, respecto de la misma circunstancia o del mismo hecho, \textbf{inhibe la instancia civil}. El proceso penal sigue adelante (probablemente se lo termine declarando culpable), pero la declaración es una forma de obtener algún beneficio o ventaja.

\section{Indictment, information y grand jury}

\begin{definition}{\textit{Indictment}}
Escrito que recibe el nombre de \textit{indictment} cuando hay un \textit{grand jury}. Es la recomendación u orden que se da al juez de «seguir adelante con la acusación». Cuando se habla de \textit{indicted}, se presupone un gran jurado, porque se trata de un delito grave al que le puede corresponder la pena de muerte.
\end{definition}

\begin{definition}{\textit{Information}}
Escrito que corresponde cuando \textbf{no hay} gran jurado. Equivale a nuestro escrito o auto de elevación a juicio, en el que ya se analizaron los indicios probatorios y hay suficiente para procesar. No lo emite el gran jurado.
\end{definition}

% TODO: contenido incompleto o ambiguo en las notas originales (quién emite la \textit{information}: pasaje poco claro en el audio; y relación entre \textit{presentment}, \textit{indictment} e \textit{information}: se afirma primero que \textit{presentment} e \textit{indictment} son «lo mismo» y luego se presenta «\textit{presentment} o \textit{information}» como el escrito cuando no hay gran jurado)

\begin{note}
En las notas, la relación entre \textit{presentment}, \textit{indictment} e \textit{information} aparece planteada de manera ambigua: se dice primero que \textit{presentment} e \textit{indictment} son lo mismo y luego se menciona «\textit{presentment} o \textit{information}» como el escrito que se emite cuando no hay gran jurado.
\end{note}

El \textbf{gran jurado} no decide la culpabilidad. Aunque no esté mal hablar de «gran jurado», en realidad es un \textbf{jurado de investigación o de acusación}: un jurado que analiza los indicios probatorios.

\section{Estándares de prueba: proceso civil y proceso penal}

\begin{formula}{Estándares de prueba en el derecho estadounidense}
\textbf{Proceso civil:} \textit{preponderance of the evidence} (preponderancia de la prueba). El jurado vota según la probabilidad; no se exige convicción absoluta.

\textbf{Proceso penal:} \textit{beyond a reasonable doubt} (más allá de toda duda razonable), con \textbf{veredicto unánime}. El jurado debe expresar que no tiene dudas de la culpabilidad; de lo contrario, debe declarar \textit{not guilty}.
\end{formula}

En el proceso civil, una vez producida la prueba, la pregunta es hacia dónde se inclina más a creer el jurado (\textit{most likely}): se trata de una cuestión de probabilidades, y los miembros del jurado votan en función de qué indicios los llevan a pensar que la persona es o no responsable. En el proceso penal, el jurado debe afirmar, «de la boca para afuera» (lo ocurrido en la sala de deliberación es otra cuestión), que no tiene dudas de que la persona es culpable.

\begin{table}[H]
\centering
\begin{tabularx}{\textwidth}{P{3.2cm}LL}
\toprule
 & \textbf{Proceso civil} & \textbf{Proceso penal} \\
\midrule
\textbf{Estándar} & \textit{Preponderance of the evidence} & \textit{Beyond a reasonable doubt} \\
\textbf{Idea central} & Probabilidad: ¿hacia dónde me inclino más a creer? & Ausencia de dudas razonables sobre la culpabilidad. \\
\textbf{Cómo decide el jurado} & Vota según la probabilidad de toda la prueba (\textit{most likely}). & Veredicto unánime; sin dudas, culpable; si no, \textit{not guilty}. \\
\textbf{Convicción requerida} & No se necesita convicción absoluta. & Convicción sin dudas (al menos «de la boca para afuera»). \\
\bottomrule
\end{tabularx}
\end{table}

\section{Cuadro Cornell: proceso penal}

\begin{table}[H]
\centering
\begin{tabularx}{\textwidth}{
    >{\raggedright\arraybackslash}p{0.28\textwidth}
    >{\raggedright\arraybackslash}X
}
\toprule
\textbf{Preguntas / palabras clave} &
\textbf{Notas / respuestas} \\
\midrule

\textbf{¿Qué es el \textit{arraignment} y cuándo se puede ajustar la fianza?} &
Audiencia inicial posterior a la detención en la que se leen formalmente los cargos y se pregunta al acusado cómo se declara. Es parecida a nuestra instrucción y ocurre antes del acuerdo con el fiscal. En esa audiencia el juez puede ajustar el monto de la fianza. \\
\midrule[0.3pt]

\textbf{¿Qué es el \textit{plea bargaining}?} &
Negociación previa con el fiscal: reducción de pena u otro beneficio a cambio de que el acusado se declare culpable, por ejemplo, para obtener información sobre otros sospechosos. \\
\midrule[0.3pt]

\textbf{¿Cómo puede declararse el acusado y qué efecto tiene el \textit{nolo contendere}?} &
\textit{Guilty}, \textit{not guilty} o \textit{nolo contendere}. Este último («no voy a contradecir») funciona como admisión de culpabilidad, pero inhibe la instancia civil por el mismo hecho. \\
\midrule[0.3pt]

\textbf{¿Qué diferencia hay entre \textit{indictment} e \textit{information}? ¿Qué es el gran jurado?} &
\textit{Indictment}: escrito del \textit{grand jury} para seguir adelante con la acusación, en delitos graves. \textit{Information}: escrito equivalente cuando no hay gran jurado, similar al auto de elevación a juicio. El gran jurado es un jurado de investigación o de acusación y no decide la culpabilidad. \\
\midrule[0.3pt]

\textbf{¿Qué estándares de prueba rigen en lo civil y en lo penal?} &
Civil: \textit{preponderance of the evidence} (probabilidad; no se exige convicción absoluta). Penal: \textit{beyond a reasonable doubt}, con veredicto unánime; si hay dudas, \textit{not guilty}. \\

\midrule

\multicolumn{2}{p{\dimexpr\textwidth-2\tabcolsep\relax}}{
\textbf{Resumen:}
Tras la detención, el \textit{arraignment} informa los cargos y pregunta cómo se declara el acusado, quien ya pudo haber negociado con el fiscal (\textit{plea bargaining}). El \textit{indictment} (con gran jurado) o la \textit{information} (sin él) habilitan seguir con la acusación. El estándar es probabilístico en lo civil y de ausencia de duda razonable, con veredicto unánime, en lo penal.
} \\

\bottomrule
\end{tabularx}
\end{table}

%==============================================================
% CAPÍTULO 3
%==============================================================
\chapter[Constitución británica, contratos y duress]{Constitución británica, contratos y \textit{duress}}

\section{La Constitución no escrita del Reino Unido}

Al estudiar derecho constitucional se vieron las diferencias entre el Reino Unido y Estados Unidos, entre ellas el hecho de que la Constitución del Reino Unido es \textbf{no escrita}. Se trabajó la \textbf{pirámide invertida}, que técnicamente no es la pirámide de las fuentes del derecho, sino la de las \textbf{fuentes de la Constitución británica}, justamente porque esta no es escrita. En ese marco se analizaron diversas diferencias a nivel teórico.

\section{Elementos y clasificación de los contratos}

Los elementos del contrato son técnicamente los mismos que en nuestro sistema, pero están \textbf{ordenados o planteados de otra manera}. El texto visto hablaba de \textit{offer and acceptance}: elementos que, en nuestro derecho, no se consideran en sí mismos, sino como parte de la \textbf{voluntad}, es decir, de los requisitos que debe haber para que un contrato sea válido.

Se vio también la \textbf{clasificación de los contratos}, según los criterios siguientes.

\begin{table}[H]
\centering
\begin{tabularx}{\textwidth}{P{5.2cm}L}
\toprule
\textbf{Criterio de clasificación} & \textbf{Variantes mencionadas} \\
\midrule
Cantidad de partes que intervienen & \textit{Unilateral contract} / \textit{bilateral contract} \\
Forma en que se comunica la aceptación & Oral, escrita o implícita (por la conducta) \\
\bottomrule
\end{tabularx}
\end{table}

\section{Unilateral contract y bilateral contract}

En el \textit{bilateral contract} ambas partes se obligan a algo y hay contraprestaciones recíprocas. En el \textit{unilateral contract} existe una \textbf{promesa, un compromiso de una sola parte}, dirigido a una \textbf{persona indeterminada o incierta}. No implica una obligación para el otro.

Si la otra persona, sea quien sea, cumple la condición, se entiende que \textbf{acepta} y el contrato se va perfeccionando. Recién entonces toma vigencia la obligación de quien hizo la promesa; hasta ese momento la promesa no se hizo efectiva, porque no hubo aceptación.

\begin{example}{Recompensa por el collar o el perro perdido}
Quien ofrece una recompensa para quien encuentre su collar hace una promesa a una persona indeterminada. La obligación del que prometió toma vigencia cuando alguien cumple la condición («encontré tu collar, encontré tu perro»): hay aceptación y se perfecciona el contrato.
\end{example}

\begin{example}{El tío y el sobrino}
Un tío le promete dinero a su sobrino a cambio de que este se abstenga de una conducta (por ejemplo, de tomar). La promesa depende de una acción o conducta del otro, lo que es característico del \textit{unilateral contract}.
\end{example}

El \textit{unilateral contract} del derecho angloamericano \textbf{no es tan equivalente} a nuestro contrato unilateral. Como ejemplos de contrato unilateral en el derecho argentino se mencionaron la donación y el testamento; la docente tomó la donación. Esta no requiere necesariamente una persona indeterminada (es para alguien específico e identificado, el destinatario) y tampoco requiere una acción de parte del otro para hacerse efectiva: es eficaz desde que la realiza una persona.

En lo que sí se parecen es en las prestaciones: «bilateral» remite a dos contraprestaciones y «unilateral» a una sola. Pero en el \textit{unilateral contract} está presente la idea de «yo prometo algo a alguien, no importa a quién, y se hará eficaz siempre que ese otro haga algo». Hay una \textbf{conducta que condiciona el contrato}, lo cual no está presente en nuestro contrato unilateral: se trata de una promesa a cambio de una acción.

% TODO: contenido incompleto o ambiguo en las notas originales (las notas editoriales del documento señalan que el testamento citado por un estudiante es en rigor un acto jurídico unilateral; la docente tomó la donación como ejemplo de contrato)

\begin{table}[H]
\centering
\begin{tabularx}{\textwidth}{P{2.8cm}LL}
\toprule
 & \textbf{\textit{Unilateral contract} (\textit{common law})} & \textbf{Contrato unilateral (derecho argentino)} \\
\midrule
\textbf{Destinatario} & Persona indeterminada o incierta («no importa a quién»). & Puede ser alguien específico e identificado (por ejemplo, donación). \\
\textbf{Acción del otro} & Se hace eficaz solo si el otro cumple una conducta o acción a cambio. & No requiere una acción de parte del otro. \\
\textbf{Eficacia} & Toma vigencia con la aceptación por cumplimiento. & Eficaz desde el momento en que lo realiza una persona. \\
\textbf{Prestaciones} & Una sola promesa, condicionada a una conducta. & Una sola prestación (en esto sí se parece). \\
\bottomrule
\end{tabularx}
\end{table}

\section{Duress: coerción y coacción}

Para traducir \textit{duress} se discutió si correspondía «coerción» o «coacción». Se distinguieron así: la \textbf{coerción} es la violencia psicológica y la \textbf{coacción}, la violencia física. Por lo tanto, \textit{duress} se traduce más bien como \textbf{coerción}: actuar a través de dichos o de insistencia para generar que alguien actúe de determinada manera, lo que afecta su voluntad.

% TODO: contenido incompleto o ambiguo en las notas originales (la docente indicó no recordar la diferencia en español; la distinción coerción/coacción fue aportada por estudiantes)

\section{Cuadro Cornell: Constitución, contratos y duress}

\begin{table}[H]
\centering
\begin{tabularx}{\textwidth}{
    >{\raggedright\arraybackslash}p{0.28\textwidth}
    >{\raggedright\arraybackslash}X
}
\toprule
\textbf{Preguntas / palabras clave} &
\textbf{Notas / respuestas} \\
\midrule

\textbf{¿Qué caracteriza a la Constitución del Reino Unido?} &
No es escrita; por eso se habla de la pirámide invertida de las fuentes de la Constitución británica (no de las fuentes del derecho en general). \\
\midrule[0.3pt]

\textbf{¿Qué lugar ocupan \textit{offer and acceptance} y cómo se clasifican los contratos?} &
Son elementos que en nuestro derecho se ven como parte de la voluntad y de los requisitos de validez. Los contratos se clasifican por la cantidad de partes (unilateral o bilateral) y por la forma de comunicar la aceptación (oral, escrita o implícita). \\
\midrule[0.3pt]

\textbf{¿En qué se diferencia el \textit{unilateral contract} de nuestro contrato unilateral?} &
El \textit{unilateral contract} es una promesa a persona indeterminada que se vuelve eficaz cuando otro cumple una conducta (recompensa por el collar). Nuestro contrato unilateral puede dirigirse a alguien identificado (donación) y no requiere acción del otro. \\
\midrule[0.3pt]

\textbf{¿Cómo se traduce \textit{duress}?} &
Como coerción: presión psicológica que induce a actuar de cierta manera y afecta la voluntad; se distingue de la coacción, asociada a la violencia física. \\

\midrule

\multicolumn{2}{p{\dimexpr\textwidth-2\tabcolsep\relax}}{
\textbf{Resumen:}
La Constitución británica no es escrita. Los elementos del contrato son los mismos que en nuestro derecho, aunque ordenados de otro modo. El \textit{unilateral contract} se perfecciona cuando otro cumple la conducta prometida y difiere de nuestro contrato unilateral. \textit{Duress} se traduce como coerción.
} \\

\bottomrule
\end{tabularx}
\end{table}

%==============================================================
% CAPÍTULO 4
%==============================================================
\chapter[El jurado]{El jurado: selección, recusaciones y garantías}

\section{Selección del jurado: azar y voir dire}

Tanto en el proceso civil como en el penal, la \textbf{primera selección} de personas es \textbf{aleatoria}, al azar: se realiza a partir del listado del registro de votantes o de la nómina de licencias de conducir. Esa primera convocatoria forma el \textit{jury panel}.

Después viene el proceso de selección, que se hace a través de \textbf{preguntas de los abogados}, con las que se trata de indagar e identificar si hay algún prejuicio o sesgo de parte de las personas convocadas, porque cada parte intenta conformar un jurado que le convenga más. Ese proceso se llama \textit{voir dire}. Por eso, cuando se dice que el jurado es aleatorio, ello es cierto de la primera selección de personas, pero después, de alguna manera, se elige a quiénes lo integrarán.

\begin{definition}{\textit{Voir dire}}
Proceso de selección de los miembros del jurado mediante preguntas de los abogados al panel, para detectar prejuicios o sesgos. En el curso del proceso algunos integrantes son excluidos.
\end{definition}

\subsection{Jury panel, jury box, jury y juror}

\begin{table}[H]
\centering
\begin{tabularx}{\textwidth}{P{3.4cm}L}
\toprule
\textbf{Término} & \textbf{Qué designa} \\
\midrule
\textit{Jury panel} & Grupo amplio de personas convocadas por sorteo: la selección preliminar («ese malón al que convocaron»). Si se necesitan 12, se llama a veinte y pico. \\
\textit{Voir dire} & Proceso de selección: preguntas al panel y exclusión de algunos de sus miembros. \\
\textit{Jury} / \textit{jury box} & Quienes quedan tras depurar el panel y efectivamente se sientan en el estrado (\textit{box}): el colectivo definitivo (12, más suplentes). \\
\textit{Juror} & Cada integrante individual del jurado definitivo (\textit{juror number one}, \textit{juror number two}\ldots). \\
\bottomrule
\end{tabularx}
\end{table}

\begin{important}
En español «jurado» designa tanto al conjunto como a cada persona. En inglés, en cambio, \textit{jury} designa al colectivo y \textit{juror} a cada integrante, y \textit{juror} se usa para el jurado definitivo, no para el panel.
\end{important}

\section{Excuse y challenge}

Los \textit{challenges} son las recusaciones, y pueden ser \textbf{con causa o sin causa}. Se distingue de la \textit{excuse} según quién plantea la objeción.

\begin{table}[H]
\centering
\begin{tabularx}{\textwidth}{P{3.4cm}P{4.2cm}L}
\toprule
\textbf{Figura} & \textbf{Quién la plantea} & \textbf{Motivo} \\
\midrule
\textit{Excuse} & La persona misma (quien fue convocada como jurado) & «No puedo porque tengo que hacer tal cosa, porque soy pariente, porque estoy involucrado.» \\
\textit{Challenge} (recusación) & Otro: una de las partes, el abogado & «Tal persona no puede desempeñarse como jurado porque\ldots»; puede ser con causa o sin causa. \\
\bottomrule
\end{tabularx}
\end{table}

\subsection{¿Quién excusa al jurado?}

La excusación parte del propio interesado (el jurado convocado), pero siempre la resuelve el \textbf{juez}: el jurado la plantea y el juez ve si corresponde. Además, el juez puede excusar a alguien apenas lo recibe en el tribunal. Ejemplo: un jurado que dice ser la única persona a cargo de un familiar enfermo; el juez lo excusa.

% TODO: contenido incompleto o ambiguo en las notas originales (la docente afirma primero que la excusación la hace solo el interesado y luego precisa que también el juez puede excusar y que siempre la resuelve el juez)

\section{Peremptory challenge y challenge for cause}

\subsection{Peremptory challenge: sin expresión de causa}

En castellano «perentorio» se usa para los plazos: algo perentorio no admite discusión. En el \textit{peremptory challenge} se aplica la misma idea respecto de la justificación: el abogado de la parte dice «yo no quiero al jurado número tal» y, si se le pregunta por qué, responde que no tiene que justificarlo. Por eso estos desafíos son \textbf{limitados}: no se pueden hacer muchas veces.

Aquí aparece el problema de la \textbf{discriminación}, que no debe existir en la conformación del jurado. Si hay, por ejemplo, cinco recusaciones sin causa y se filtra discriminación de modo sistemático (excluir a todas las personas de cierta religión o a todas las mujeres), seguramente la contraparte objetará o planteará el punto, o el juez, de oficio, lo señalará. Hay allí una posibilidad de abuso, porque la recusación es sin expresión de causa.

\subsection{Challenge for cause: con expresión de causa}

En el \textit{challenge for cause} se expresa la causa: por ejemplo, que la persona tuvo participación en el hecho, que tiene un vínculo de amistad, o que atravesó una situación similar y se sabe que tiene cierta forma de pensar. Estos desafíos son \textbf{ilimitados}: habrá tantos como pueda justificar la parte, y el juez verá si hace lugar o no.

\begin{table}[H]
\centering
\begin{tabularx}{\textwidth}{P{3.2cm}LL}
\toprule
 & \textbf{\textit{Peremptory challenge}} & \textbf{\textit{Challenge for cause}} \\
\midrule
\textbf{¿Se expresa la causa?} & No: «no tengo que justificar por qué». & Sí: se expresa y se justifica. \\
\textbf{Cantidad} & Limitados. & Ilimitados: tantos como pueda justificar. \\
\textbf{Ejemplos} & «Yo no quiero al jurado número tanto». & Participó en el hecho, vínculo de amistad, situación similar y forma de pensar. \\
\textbf{Control} & Si se vuelve sistemático (excluir a todas las mujeres o a todas las personas de una religión), la contraparte puede objetar o el juez actuar de oficio. & El juez decide si hace lugar o no. \\
\bottomrule
\end{tabularx}
\end{table}

\section{Garantías de quienes integran el jurado}

La convocatoria a los miembros del jurado es aleatoria, y existen \textbf{garantías} para que quienes son convocados no pierdan su trabajo.

\begin{important}
\textbf{Garantías para los miembros del jurado (según se explicó en clase).}
El empleador no puede despedir a quien es convocado y debe conservarle el puesto de trabajo. No está obligado a pagarle el sueldo si la justicia o el Estado abona un estipendio o viático (comida, transporte y demás) mientras dura el proceso. Está prohibido discriminar en la selección del jurado, al menos de manera manifiesta.
\end{important}

Se señaló que ello depende de las circunstancias y que hay situaciones, como la de los \textbf{autónomos} o independientes, que afectan o ponen en tela de juicio el buen funcionamiento del juicio por jurado. Además, la forma en que se hacen las recusaciones obliga a analizar si realmente se cumple el principio de no discriminación.

\section{Trier of fact: jury trial y bench trial}

Los miembros del jurado son el \textit{trier of fact}: el jurado \textbf{resuelve sobre los hechos} y el juez \textbf{sobre el derecho}. Cuando no hay jurado y todo lo lleva adelante el juez, el juicio es un \textit{bench trial}, más parecido a nuestro sistema; en ese caso el juez es también el \textit{trier of fact}. Con jurado, el juicio es un \textit{jury trial}.

No existe una palabra en español para traducir \textit{trier}: es quien resuelve o analiza los hechos.

\begin{table}[H]
\centering
\begin{tabularx}{\textwidth}{P{4.4cm}LL}
\toprule
 & \textbf{\textit{Jury trial}} & \textbf{\textit{Bench trial}} \\
\midrule
\textbf{Quién resuelve los hechos (\textit{trier of fact})} & El jurado & El juez \\
\textbf{Quién resuelve el derecho} & El juez & El juez \\
\textbf{Comparación con Argentina} & --- & Más parecido a nuestro sistema \\
\bottomrule
\end{tabularx}
\end{table}

\section{Cuadro Cornell: el jurado}

\begin{table}[H]
\centering
\begin{tabularx}{\textwidth}{
    >{\raggedright\arraybackslash}p{0.28\textwidth}
    >{\raggedright\arraybackslash}X
}
\toprule
\textbf{Preguntas / palabras clave} &
\textbf{Notas / respuestas} \\
\midrule

\textbf{¿Cómo se distinguen \textit{jury panel}, \textit{voir dire}, \textit{jury} y \textit{juror}?} &
\textit{Jury panel}: grupo convocado por sorteo. \textit{Voir dire}: selección con preguntas de los abogados. \textit{Jury} / \textit{jury box}: el jurado definitivo (12, más suplentes). \textit{Juror}: cada integrante. \\
\midrule[0.3pt]

\textbf{¿Qué diferencia hay entre \textit{excuse} y \textit{challenge}?} &
\textit{Excuse}: la plantea la persona convocada (no puede, es pariente, está involucrada); la resuelve el juez. \textit{Challenge} (recusación): la plantean las partes contra un integrante del jurado. \\
\midrule[0.3pt]

\textbf{¿Cuándo hay que expresar la causa de la recusación?} &
\textit{Peremptory challenge}: sin expresar causa y en cantidad limitada (si es sistemático puede haber discriminación). \textit{Challenge for cause}: se expresa y justifica la causa; son ilimitados y el juez decide si hace lugar. \\
\midrule[0.3pt]

\textbf{¿Qué garantías tienen quienes son convocados como jurado?} &
No pueden ser despedidos y conservan el puesto; el empleador no debe pagar el sueldo si el Estado paga un estipendio o viático; está prohibido discriminar en la selección. \\
\midrule[0.3pt]

\textbf{¿Quién es el \textit{trier of fact}?} &
Quien resuelve los hechos: el jurado en el \textit{jury trial} y el juez en el \textit{bench trial}. El juez resuelve el derecho. No tiene traducción directa al español. \\

\midrule

\multicolumn{2}{p{\dimexpr\textwidth-2\tabcolsep\relax}}{
\textbf{Resumen:}
El jurado se forma a partir de un panel convocado al azar, depurado mediante el \textit{voir dire}, las excusas y las recusaciones (\textit{peremptory} y \textit{for cause}). Quienes lo integran tienen garantías laborales y de no discriminación. El jurado resuelve los hechos (\textit{trier of fact}) y el juez el derecho; sin jurado, el juez resuelve ambos (\textit{bench trial}).
} \\

\bottomrule
\end{tabularx}
\end{table}

%==============================================================
% CAPÍTULO 5
%==============================================================
\chapter[Motions y equitable remedies]{\textit{Motions} y \textit{equitable remedies}}

\section{Motion: significado y traducción}

\begin{definition}{\textit{Motion}}
En sentido amplio, un pedido, una solicitud o un requerimiento que una de las partes hace al juez para que resuelva algo, en cualquier instancia del proceso.
\end{definition}

\textit{Motion} se traduce como pedido, solicitud o requerimiento, según el caso. Si lo que se solicita es, por ejemplo, una impugnación, la \textit{motion} puede traducirse en consecuencia, pero a grandes rasgos es esto: una parte le pide algo al juez.

\begin{important}
\textit{Motion} \textbf{no} se traduce como «moción»: en español, moción es una propuesta que otros votan o aprueban por sí o por no, no un pedido a un juez para que resuelva. Es un falso amigo.
\end{important}

\section{Momentos procesales: pretrial, trial y post-trial}

Las \textit{motions} pueden presentarse en tres momentos: \textit{pretrial}, \textit{trial} y \textit{post-trial}.

\subsection{Pretrial motions}

Siguen siendo pedidos al juez, pero se formulan \textbf{mucho antes de la etapa probatoria}, cuando todavía se está instalando la causa (\textit{early stages}), es decir, cuando se traba la litis y se está en la etapa de la demanda. Se las puede asimilar a nuestras \textbf{excepciones previas}. Ejemplos:

\begin{itemize}
    \item \textit{Motion to dismiss}: pedido de desestimar la demanda.
    \item \textit{Summary judgment}: también es una \textit{motion}, un pedido al juez para que declare la causa de puro derecho, porque no hay hechos controvertidos (por ejemplo, porque la parte demandada se allanó a los hechos y reconoció que son ciertos). El juez resuelve entonces de puro derecho.
\end{itemize}

Alcanza con decir que son «pedidos o solicitudes al juez para que resuelva en las etapas iniciales del proceso».

\subsection{Trial motions}

Tienen que ver con \textbf{cuestiones probatorias}: ya se está en la etapa probatoria y los pedidos refieren a la prueba misma («esta prueba no es admisible», «esto sí», «tal testigo está mintiendo»). Se mencionaron como posibles traducciones impugnaciones, admisiones y denegaciones.

\subsection{Post-trial motions}

Una vez concluida la etapa probatoria, cuando ya hubo veredicto y se dictó sentencia, comprenden todo lo que las partes puedan pedirle al juez para que \textbf{reevalúe o revea} algo: apelaciones y recursos (de queja, de revisión, de apelación).

\begin{table}[H]
\centering
\begin{tabularx}{\textwidth}{P{2cm}LLL}
\toprule
\textbf{Momento} & \textbf{Qué se pide al juez} & \textbf{Cómo se puede traducir} & \textbf{Ejemplos} \\
\midrule
\textit{Pretrial} & Que resuelva algo mucho antes de la etapa probatoria (etapas iniciales). & Pedido o solicitud al juez; asimilable a las excepciones previas. & \textit{Motion to dismiss}; \textit{summary judgment}. \\
\textit{Trial} & Cuestiones de la etapa probatoria y de la prueba misma. & Impugnaciones, admisiones, denegaciones. & Una prueba no es admisible; un testigo está mintiendo. \\
\textit{Post-trial} & Que el juez revea o reevalúe algo tras el veredicto y la sentencia. & Apelaciones, recursos (queja, revisión). & Recursos contra la sentencia. \\
\bottomrule
\end{tabularx}
\end{table}

Según el momento, quien traduce puede elegir cómo hacerlo; si se trata de la primera etapa puede hablar de excepciones si quiere asimilarla, pero alcanza con decir «pedido o solicitud al juez para que resuelva algo».

\section{Equitable remedies}

Son tres: \textit{injunction}, \textit{specific performance} y \textit{rescission}.

\begin{table}[H]
\centering
\begin{tabularx}{\textwidth}{P{3.6cm}LL}
\toprule
\textbf{\textit{Equitable remedy}} & \textbf{Qué es} & \textbf{Alcance} \\
\midrule
\textit{Injunction} & Orden de no hacer. & Causas civiles, no solo contratos. \\
\textit{Specific performance} & Orden de hacer lo que estaba prometido u obligado. & Causas civiles, no solo contratos. \\
\textit{Rescission} & Resolución: dejar sin efecto el contrato y retrotraer los efectos. & Solo contratos. \\
\bottomrule
\end{tabularx}
\end{table}

\section{Origen histórico y problema de traducción}

El adjetivo \textit{equitable} \textbf{no} viene de que estos remedios sean «equitativos» o justos en el sentido de \textit{fair}: la indemnización por daños y perjuicios también es justa, llegado el caso, si el monto está bien determinado.

Proviene del contexto histórico en que los poderes delegados de la corona o del rey se trasladaron a \textbf{jueces itinerantes}. En un momento, la aplicación de la analogía y del precedente judicial tomó impulso: se produjeron variaciones y se encontraron alternativas de solución a conflictos que antes no las tenían. Antes de estos \textit{remedies} existían los \textbf{\textit{legal remedies}}, que eran ni más ni menos que la indemnización por daños y perjuicios (\textit{damages}); era lo único que había originalmente. Como había situaciones que no se podían resolver de esa manera, surgieron los \textit{equitable remedies}, un sistema paralelo o nueva corriente.

Hoy un juez puede ordenar cualquiera de estas medidas y ya no hay diferencia ni tribunales diferenciados. Pero cuando en un texto aparece \textit{equitable}, hace mucho ruido: los traductores automáticos dicen «equitativo», lo cual induce a error. La diferencia ya tiene poca razón de ser, pero se sigue la terminología y presenta problemas al mundo, ya que se supone que la justicia es siempre equitativa.

\begin{table}[H]
\centering
\begin{tabularx}{\textwidth}{P{2.4cm}LL}
\toprule
 & \textbf{\textit{Legal remedies}} & \textbf{\textit{Equitable remedies}} \\
\midrule
\textbf{Contenido} & Indemnización por daños y perjuicios (\textit{damages}). & \textit{Injunction}, \textit{specific performance}, \textit{rescission}. \\
\textbf{Origen} & Era lo único que había originalmente. & Surgen de jueces itinerantes que, con la analogía y el precedente, buscan alternativas para conflictos sin solución. \\
\textbf{Hoy} & Un mismo juez puede ordenar cualquiera de las medidas. & Ya no hay tribunales diferenciados; la diferencia casi no tiene razón de ser, pero la terminología se mantiene. \\
\bottomrule
\end{tabularx}
\end{table}

\section{Cuadro Cornell: motions y equitable remedies}

\begin{table}[H]
\centering
\begin{tabularx}{\textwidth}{
    >{\raggedright\arraybackslash}p{0.28\textwidth}
    >{\raggedright\arraybackslash}X
}
\toprule
\textbf{Preguntas / palabras clave} &
\textbf{Notas / respuestas} \\
\midrule

\textbf{¿Cómo se traduce \textit{motion}?} &
Pedido, solicitud o requerimiento de una parte al juez para que resuelva algo; no «moción», porque moción es una propuesta que otros votan. \\
\midrule[0.3pt]

\textbf{¿En qué momentos se presentan las \textit{motions}?} &
\textit{Pretrial} (desestimar la demanda, \textit{summary judgment}; asimilables a excepciones previas), \textit{trial} (cuestiones probatorias) y \textit{post-trial} (apelaciones y recursos tras el veredicto y la sentencia). \\
\midrule[0.3pt]

\textbf{¿Qué son los \textit{equitable remedies}?} &
\textit{Injunction} (orden de no hacer), \textit{specific performance} (orden de hacer lo prometido) y \textit{rescission} (resolución, solo contratos). \\
\midrule[0.3pt]

\textbf{¿Por qué \textit{equitable} no significa «equitativo»?} &
Porque el término es histórico: designa un sistema paralelo de remedios surgido frente a los \textit{legal remedies} (\textit{damages}), que eran los únicos originalmente. La indemnización por daños también es justa. \\

\midrule

\multicolumn{2}{p{\dimexpr\textwidth-2\tabcolsep\relax}}{
\textbf{Resumen:}
\textit{Motion} es un pedido al juez, no una «moción», y puede presentarse antes, durante o después del juicio. Los \textit{equitable remedies} (\textit{injunction}, \textit{specific performance}, \textit{rescission}) son un resto histórico frente a los \textit{legal remedies}, y \textit{equitable} no significa «equitativo».
} \\

\bottomrule
\end{tabularx}
\end{table}

%==============================================================
% CAPÍTULO 6
%==============================================================
\chapter[Discovery y depositions]{\textit{Discovery}, \textit{depositions} y prueba nueva}

\section{Discovery: una etapa preliminar híbrida}

Al estudiar el proceso civil se fue encajando cada etapa en alguna de las nuestras. El \textit{discovery} se corresponde con la \textbf{etapa de ofrecimiento de prueba}, pero el momento no coincide: nosotros ofrecemos las pruebas junto con la demanda, de modo que el \textit{discovery} es \textbf{más amplio} que el ofrecimiento. Es una instancia en la que los abogados «tantean» y se intercambian la prueba.

\begin{definition}{\textit{Discovery}}
Etapa preliminar de intercambio de pruebas e información probatoria entre las partes. A veces se habla de «revelación», en el sentido de revelar a la contraparte cuáles serán las pruebas que se van a producir.
\end{definition}

\subsection{Supervisión del juez}

El \textit{discovery} no se realiza en presencia del juez, pero el juez \textbf{está}: un taquígrafo registra y transcribe las actuaciones, las partes pueden acudir al juez para informarle lo ocurrido y el juez las convoca cada tanto para preguntar cómo van, si falta mucho, si se puede avanzar o si hay acuerdo. No está presente, pero participa.

\begin{important}
El \textit{discovery} \textbf{no es extrajudicial}: es un error muy común decirlo. Es verdad que no sucede en la sala del tribunal, pero ya hay un proceso judicial en curso. Lo que es es \textbf{preliminar}: un intercambio de pruebas y de información probatoria.
\end{important}

\subsection{Acuerdo (settlement)}

Al terminar la etapa puede ocurrir que una de las partes se baje del reclamo y se llegue a un \textit{settlement}, porque advirtió que no va a ganar y que seguir con todo el proceso es perder tiempo. En ese caso se informa al juez que concluye el \textit{discovery} y que hay acuerdo; no hace falta jurado, porque no hay etapa probatoria ni \textit{trial}. El juez instrumenta las cuestiones de derecho que haga falta y dicta sentencia, sin intervención de jurado, de modo que el proceso se acorta y se abrevia considerablemente.

\section{Depositions}

Las \textit{depositions} son declaraciones que se toman \textbf{entre los abogados}, \textbf{bajo juramento}. Muchas veces se filman y se transcriben, aunque parece que se transcribe menos que antes. Si se detecta que la persona miente, las partes pueden plantearlo al juez.

No son todavía la prueba testimonial: son \textbf{anteriores} a ella, ya que durante el \textit{discovery} todavía no es prueba, porque aún no está el jurado ni se llegó a la etapa probatoria. Pueden verse como un \textbf{anticipo} de la prueba testimonial que permite ir más preparado.

La etapa es un \textbf{híbrido}: es privada solo en el sentido de que no ocurre en la sala del tribunal ante el juez, pero la supervisión del juez está siempre y las partes pueden acudir a él en cualquier momento («señor juez, pasó tal cosa: fulano se levantó en medio de la audiencia y se fue, insultó, no declaró, mintió, no lo pudimos convocar, no lo encontramos»). No es prejudicial ni extrajudicial.

\begin{table}[H]
\centering
\begin{tabularx}{\textwidth}{P{3.2cm}L}
\toprule
\textbf{Aspecto} & \textbf{\textit{Depositions}} \\
\midrule
\textbf{Quiénes participan} & Los abogados entre sí. \\
\textbf{Carácter} & Bajo juramento; a veces filmadas y transcriptas. \\
\textbf{Valor} & Anticipo de la prueba testimonial; todavía no es prueba. \\
\textbf{Control} & Si se detecta que alguien miente, se acude al juez. \\
\bottomrule
\end{tabularx}
\end{table}

\section{Cierre del discovery y prueba nueva}

\subsection{Cómo se cierra la etapa}

El juez va preguntando si terminaron con todos los testigos, si ya recibieron todas las \textit{depositions} y si se puede cerrar la etapa y avanzar. Al cerrarla, hay dos posibilidades: si \textbf{hay acuerdo}, se cierra y el juez resuelve las cuestiones de derecho, sin jurado. Si \textbf{no hay acuerdo}, la causa se abre a prueba y entonces se inicia la convocatoria de los miembros del jurado.

\subsection{La prueba no ofrecida: preclusión y excepciones}

La etapa es \textbf{preclusiva en teoría}: el principio es que todo lo que no se propone o no se revela en el \textit{discovery} no se puede llevar después como prueba. Pero si alguna de las partes puede justificar por qué no presentó, no mencionó o no intercambió determinada prueba, el juez lo analizará. En el momento del \textit{trial} puede aparecer una prueba nueva que no se presentó en el \textit{discovery}: se admitirá, probablemente, si se justifica como un \textbf{hecho nuevo} o prueba nueva, pero habrá una objeción de la otra parte (``esa prueba yo no la conozco'', «¿por qué apareció ahora?»).

\begin{example}{Prueba nueva en el trial}
Un testigo que viene del extranjero y que la parte no conocía o creyó que no iba a estar: la parte deberá justificar ante el juez por qué no lo presentó antes, la contraparte objetará y el juez decidirá si la admite.
\end{example}

Esto se vincula con la \textbf{estrategia} del abogado: a veces se puede usar un hecho nuevo como estrategia (``me lo guardé, me hice el tonto, porque no te lo quería anticipar''). La docente señaló que este es otro tema bueno para desarrollar: la anticipación de pruebas. También puede ocurrir que un testigo o una prueba que se consideró no relevante resulte después de utilidad en el proceso.

Por eso los abogados tratan de ser \textbf{lo más exhaustivos posible} en el intercambio: como lo que no se menciona en el \textit{discovery} técnicamente no puede presentarse después, ofrecen mucha prueba. Después, en el proceso en sí, no tiene tanta relevancia, pero ya está mencionada e introducida, y se la puede usar a favor o en contra, según el caso.

\begin{important}
\textbf{Regla general y estrategia:} lo que no se propone o no se revela en el \textit{discovery}, en principio no puede presentarse después (efecto preclusivo). Puede admitirse si la parte justifica un hecho nuevo, pero habrá objeción de la otra parte.
\end{important}

\section{Cuadro Cornell: discovery}

\begin{table}[H]
\centering
\begin{tabularx}{\textwidth}{
    >{\raggedright\arraybackslash}p{0.28\textwidth}
    >{\raggedright\arraybackslash}X
}
\toprule
\textbf{Preguntas / palabras clave} &
\textbf{Notas / respuestas} \\
\midrule

\textbf{¿Qué es el \textit{discovery} y por qué no es extrajudicial?} &
Etapa preliminar híbrida de intercambio de información probatoria, más amplia que nuestro ofrecimiento de prueba. Hay un proceso judicial en curso y el juez supervisa y puede ser consultado en cualquier momento, aunque no se desarrolle en la sala del tribunal. \\
\midrule[0.3pt]

\textbf{¿Qué pasa si hay \textit{settlement}?} &
Se informa al juez, que resuelve las cuestiones de derecho y dicta sentencia sin jurado: no hay etapa probatoria ni \textit{trial} y el proceso se abrevia. \\
\midrule[0.3pt]

\textbf{¿Qué son las \textit{depositions} y qué valor tienen?} &
Declaraciones bajo juramento entre los abogados, a veces filmadas o transcriptas. Son un anticipo de la prueba testimonial, todavía no prueba; si se detecta que alguien miente, se acude al juez. \\
\midrule[0.3pt]

\textbf{¿Qué pasa con la prueba no ofrecida en el \textit{discovery}?} &
En principio no puede presentarse después (preclusión), salvo que se justifique un hecho nuevo y el juez lo admita pese a la objeción de la otra parte. Los abogados buscan ser exhaustivos en el intercambio. \\

\midrule

\multicolumn{2}{p{\dimexpr\textwidth-2\tabcolsep\relax}}{
\textbf{Resumen:}
El \textit{discovery} es una etapa judicial preliminar, supervisada por el juez, de intercambio de prueba y de \textit{depositions} bajo juramento. Puede terminar en un \textit{settlement}, que evita el jurado, o en la apertura a prueba. Lo no revelado en esta etapa, en principio, precluye, salvo hecho nuevo justificado.
} \\

\bottomrule
\end{tabularx}
\end{table}

\end{document}
